\documentclass[twocolumn]{article}
\usepackage[utf8]{inputenc}
\usepackage{xcolor}
\usepackage{fancyhdr}
\usepackage{amsmath}

\title{Preprint Científico: Modelo de Partículas y Cómputo Analítico}
\author{Núcleo Inteligente de Disboard-NGNL}
\date{\today}

% --- SUNI_STYLIST WILL INJECT HERE ---
\begin{document}
\maketitle

\begin{abstract}
Este preprint académico expone el modelado analítico de partículas cuánticas y de física matemática implementado bajo la orquestación multiagente 『 』.
\end{abstract}

\section{Introducción}
El modelado cuántico simbólico en el codespace Suniaster aprovecha la inferencia de bajo consumo de tokens.

\section{Metodología}
La suite de cálculo Shiro-Kurami define las restricciones matemáticas:
\begin{equation}
i\hbar \frac{\partial}{\partial t} \Psi = \hat{H} \Psi
\end{equation}

\section{Resultados}
Los datos experimentales provistos por Izuna Hatsuse confirman el acoplamiento óptimo.

\bibliographystyle{plain}
\begin{thebibliography}{1}
\bibitem{ngnl} Tet (2026). \emph{Las Diez Reglas del Juego de Disboard}. Editorial Suniaster.
\end{thebibliography}

\end{document}
